% ======================================================
%  Sekcja 5 — Mieszane Wyzwania
%  5 zadań · prędkość, procenty, geometria, finanse, praca
% ======================================================

% -- 1 --
\ExerciseSlide[\CategoryBadge[MixColor!20]{Zakupy}]{1}{30 s}{%
    Kawa kosztuje 12~z\l{}, p\l{}acisz banknotem 50~z\l{}.\par
    Ile dostaniesz reszty?
  }
\AnswerSlide[\CategoryBadge[MixColor!20]{Zakupy}]{38~z\l{}}{$50 - 12 = 38$}

% -- 2 --
\ExerciseSlide[\CategoryBadge[MixColor!20]{Czas}]{1}{60 s}{%
    Film zaczyna si\k{e} o 18:45 i trwa 95~minut.\par
    O kt\'{o}rej si\k{e} ko\'nczy?
  }
\AnswerSlide[\CategoryBadge[MixColor!20]{Czas}]{20:20}{95~min to 1~h 35~min.}

% -- 3 --
\ExerciseSlide[\CategoryBadge[MixColor!20]{Pr\k{e}dko\'s\'c}]
  {2}{2 min}{%
    Poci\k{a}g jedzie 120~km w 1{,}5~h.\par
    Jaka jest jego \'{s}rednia pr\k{e}dko\'s\'c?
  }
\AnswerSlide[\CategoryBadge[MixColor!20]{Pr\k{e}dko\'s\'c}]{80~km/h}{$120 \div 1{,}5 = 80$}

% -- 4 --
\ExerciseSlide[\CategoryBadge[MixColor!20]{Procenty}]
  {2}{2 min}{%
    Pizza ma 8 kawa\l{}k\'{o}w, zjadasz 3.\par
    Jaki procent zostaje?
  }
\AnswerSlide[\CategoryBadge[MixColor!20]{Procenty}]{62{,}5\%}{$\tfrac{5}{8} \times 100\% = 62{,}5\%$}

% -- 5 --
\ExerciseSlide[\CategoryBadge[MixColor!20]{Geometria}]
  {3}{3 min}{%
    Pok\'{o}j $5\,\text{m}\times 4\,\text{m}$.\par
    Kafelki $50\,\text{cm}\times 50\,\text{cm}$.\par
    Ile kafelk\'{o}w?
  }
\AnswerSlide[\CategoryBadge[MixColor!20]{Geometria}]{80}{$10 \times 8 = 80$}

% -- 6 --
\ExerciseSlide[\CategoryBadge[MixColor!20]{Finanse}]
  {3}{3 min}{%
    1000~z\l{} przy 5\% rocznie.\par
    Procent sk\l{}adany. 2~lata.\par
    Ile masz?
  }
\AnswerSlide[\CategoryBadge[MixColor!20]{Finanse}]{1102{,}50~z\l{}}{$1000 \times 1{,}05^2$}

% -- 7 --
\ExerciseSlide[\CategoryBadge[MixColor!20]{Praca}]
  {3}{5 min}{%
    3 robotnik\'{o}w robi prac\k{e} w 6 dni.\par
    Ile dni zajmie 2 robotnikom?
  }
\AnswerSlide[\CategoryBadge[MixColor!20]{Praca}]{9 dni}{$3 \times 6 = 18$ robotnikodni $\div$ $2 = 9$.}

% -- 8 --
\ExerciseSlide[\CategoryBadge[MixColor!20]{Proporcje}]{2}{90 s}{%
    Przepis dla 4~os\'{o}b wymaga 300~g m\k{a}ki.\par
    Ile m\k{a}ki na 6~os\'{o}b?
  }
\AnswerSlide[\CategoryBadge[MixColor!20]{Proporcje}]{450~g}{$300 \div 4 = 75$~g na osob\k{e}.}
